\section{Problématique}

\begin{frame}{Du terrain au problème}
  \footnotesize
  \begin{columns}[T,onlytextwidth]
    \begin{column}{0.50\textwidth}
      \begin{block}{Le terrain}
        \begin{itemize}
          \item 582 patients, 280 vacations sur 40 jours.
          \item Trois ressources couplées : bloc, lits, ambulatoire.
          \item Quatre aléas : urgence, annulation, fermeture, retard.
        \end{itemize}
      \end{block}

      \uncover<2->{%
      \begin{alertblock}{Le nombre d'affectations}
        \[
          \underbrace{280^{582}}_{\text{affectations possibles}}
        \]
        \footnotesize Les compatibilités réduisent l'espace, sans le rendre
        traitable.
      \end{alertblock}}
    \end{column}

    \begin{column}{0.48\textwidth}
      \centering
      \begin{tikzpicture}[
        c/.style={draw=McmTerracotta,rounded corners=3pt,fill=white,
                  text width=4.5cm,minimum height=1.10cm,align=center,
                  font=\footnotesize,inner sep=3pt},
        f/.style={-{Latex[length=2.5mm]},thick,draw=McmTerracotta}]
        \node[c] (a) at (0,0) {\textbf{Terrain}\\3 ressources, 4 aléas};
        \node[c] (b) at (0,-1.7) {\textbf{Modèle}\\affectation sous contraintes};
        \node<2->[c,draw=McmAmber,fill=McmAmber!8] (c) at (0,-3.3)
          {\textbf{Problème NP-difficile}\\qualité en temps borné};
        \draw[f] (a) -- (b);
        \draw<2->[f,draw=McmAmber] (b) -- (c);
      \end{tikzpicture}
      \uncover<3->{\par\vspace{0.15cm}\footnotesize Dix méthodes, un même
        budget de temps.}
    \end{column}
  \end{columns}
  \source{Charges et protocoles : RAPPORT.md.}
\end{frame}

\begin{frame}{Les contraintes du terrain}
  \scriptsize
  \begin{columns}[T,onlytextwidth]
    \begin{column}{0.49\textwidth}
      \begin{block}{À l'exécution}
        \begin{itemize}
          \item \textbf{Gel du passé} : un acte commencé ou programmé avant le
            jour courant n'est jamais déplacé.
          \item \textbf{Urgences} : insertion dans la fenêtre clinique, pas dès
            qu'une place reste.
          \item \textbf{Moindre perturbation} : déplacer le moins possible les
            patients déjà prévenus.
        \end{itemize}
      \end{block}
    \end{column}

    \begin{column}{0.49\textwidth}
      \begin{alertblock}{CE et RGPD}
        \begin{itemize}
          \item \textbf{CE} : logiciel d'aide à la décision à qualifier au
            titre du règlement (UE) 2017/745.
          \item \textbf{RGPD} : minimisation, pseudonymisation, base légale,
            hébergement de données de santé.
          \item \textbf{Audit} : qui a décidé quoi, quand et pourquoi.
        \end{itemize}
      \end{alertblock}
    \end{column}
  \end{columns}

  \uncover<2->{\par\vspace{0.12cm}\footnotesize Le gel, la traçabilité et
    l'explicabilité contraignent le voisinage, le score et l'interface.}
  \source{Sources : règlement (UE) 2017/745 ; RGPD (UE) 2016/679.}
\end{frame}
